%%%%%%%%%%%%%%%%%%%%%%%%%%%%%%%%%%%%%%%%%%%%%%%%%%%%%%%%%%%%%%%%%%%%%%%%%%%%%%%%%%%%%%%%%%%%%%%%%%%
%
% Primordial Machine's Math Matrices Library
% Copyright (c) 2017-2019 Michael Heilmann
%
% This software is provided 'as-is', without any express or implied warranty.
% In no event will the authors be held liable for any damages arising from the
% use of this software.
%
% Permission is granted to anyone to use this software for any purpose,
% including commercial applications, and to alter it and redistribute it
% freely, subject to the following restrictions:
%
% 1. The origin of this software must not be misrepresented;
%    you must not claim that you wrote the original software.
%    If you use this software in a product, an acknowledgment
%    in the product documentation would be appreciated but is not required.
%
% 2. Altered source versions must be plainly marked as such,
%    and must not be misrepresented as being the original software.
%
% 3. This notice may not be removed or altered from any source distribution.
%
%%%%%%%%%%%%%%%%%%%%%%%%%%%%%%%%%%%%%%%%%%%%%%%%%%%%%%%%%%%%%%%%%%%%%%%%%%%%%%%%%%%%%%%%%%%%%%%%%%%

\documentclass[oneside]{book}

%%%%%%%%%%%%%%%%%%%%%%%%%%%%%%%%%%%%%%%%%%%%%%%%%%%%%%%%%%%%%%%%%%%%%%%%%%%%%%%%%%%%%%%%%%%%%%%%%%%
%
% Primordial Machine's Math Matrices Library
% Copyright (C) 2017-2019 Michael Heilmann
%
% This software is provided 'as-is', without any express or implied warranty.
% In no event will the authors be held liable for any damages arising from the
% use of this software.
%
% Permission is granted to anyone to use this software for any purpose,
% including commercial applications, and to alter it and redistribute it
% freely, subject to the following restrictions:
%
% 1. The origin of this software must not be misrepresented;
%    you must not claim that you wrote the original software.
%    If you use this software in a product, an acknowledgment
%    in the product documentation would be appreciated but is not required.
%
% 2. Altered source versions must be plainly marked as such,
%    and must not be misrepresented as being the original software.
%
% 3. This notice may not be removed or altered from any source distribution.
%
%%%%%%%%%%%%%%%%%%%%%%%%%%%%%%%%%%%%%%%%%%%%%%%%%%%%%%%%%%%%%%%%%%%%%%%%%%%%%%%%%%%%%%%%%%%%%%%%%%%

\usepackage[utf8]{inputenc}
\usepackage{textcomp} % for \textlangle and \textrangle macros
\usepackage{xcolor}
%
\usepackage{datetime}
%
\usepackage{tabularx} % for tabularx environment
\usepackage{multirow}
%
\usepackage{lmodern}
\usepackage{microtype}
%
\usepackage{graphicx}
\usepackage[colorlinks]{hyperref}
\usepackage[margin=0.75in]{geometry}
\usepackage{amsmath,amssymb,amsfonts}

\setlength{\parindent}{0pt} % Disable paragraph indention.

\makeatletter

% "(Get|Set)Author".
% The name of the author.
\def\SetAuthor#1{\gdef\@author{#1}}
\def\@author{\@latex@warning@no@line{No author given}}
\def\GetAuthor{\@author}

% "(Get|Set)Email".
% The email (address) of the author.
\def\SetEmail#1{\gdef\@email{#1}}
\def\@email{\@latex@warning@no@line{No email given}}
\def\GetEmail{\@email}

% "(Get|Set)Organization".
% The organization the library is published by.
\def\SetOrganization#1{\gdef\@organization{#1}}
\def\@organization{\@latex@warning@no@line{No organization given}}
\def\GetOrganization{\@organization}

% "(Get|Set)LibraryName".
% The name of the library.
\def\SetLibraryName#1{\gdef\@libraryName{#1}}
\def\@libraryName{\@latex@warning@no@line{No library name given}}
\def\GetLibraryName{\@libraryName}

% "(Get|Set)LibraryIncludesPath".
% The includes directory path.
% Example: "primordialmachine/arithmetic-functors/$(PlatformTarget.toLower())/$(Configuration.toLower())/includes"
\def\SetLibraryIncludesDirectoryPath#1{\gdef\@libraryIncludesDirectoryPath{#1}}
\def\@libraryIncludesDirectoryPath{\@latex@warning@no@line{No library includes directory path given}}
\def\GetLibraryIncludesDirectoryPath{\@libraryIncludesDirectoryPath}

% "(Get|Set)LibraryIncludeFile".
% The include file name of the library.
% Example: "include.hpp".
\def\SetLibraryIncludeFileName#1{\gdef\@libraryIncludeFileName{#1}}
\def\@libraryIncludeFileName{\@latex@warning@no@line{No library include file name given}}
\def\GetLibraryIncludeFileName{\@libraryIncludeFileName}

% "(Get|Set)LibraryIncludeDirectiveFilePath".
% The library include directive file path.
% Example: "primordialmachine/arithmetic_functors/include.hpp".
\def\SetLibraryIncludeDirectiveFilePath#1{\gdef\@libraryIncludeDirectiveFilePath{#1}}
\def\@libraryIncludeDirectiveFilePath{\@latex@warning@no@line{No library include directive file path}}
\def\GetLibraryIncludeDirectiveFilePath{\@libraryIncludeDirectiveFilePath}

% "(Get|Set)LibraryStaticLibrariesDirectoryPath".
% The static libraries directory path.
% Example: "primordialmachine/arithmetic-functors/$(PlatformTarget.toLower())/$(Configuration.toLower())/libraries".
\def\SetLibraryStaticLibrariesDirectoryPath#1{\gdef\@libraryStaticLibrariesDirectoryPath{#1}}
\def\@libraryStaticLibrariesDirectoryPath{\@latex@warning@no@line{No library static libraries directory path given}}
\def\GetLibraryStaticLibrariesDirectoryPath{\@libraryStaticLibrariesDirectoryPath}

% "(Get|Set)LibraryStaticLibraryFileName.
% The static library file name of the library.
% Example: "arithmetic-functors.hpp".
\def\SetLibraryStaticLibraryFileName#1{\gdef\@libraryStaticLibraryFileName{#1}}
\def\@libraryStaticLibraryFileName{\@latex@warning@no@line{No library static library file name given}}
\def\GetLibraryStaticLibraryFileName{\@libraryStaticLibraryFileName}

% "(Get|Set)LibraryVersion".
% The version of the library.
\def\SetLibraryVersion#1{\gdef\@libraryVersion{#1}}
\def\@libraryVersion{\@latex@warning@no@line{No library version given}}
\def\GetLibraryVersion{\@libraryVersion}

% "(Get|Set)LibraryRepository".
% The repository (url) of the library.
\def\SetLibraryRepository#1{\gdef\@libraryRepository{#1}}
\def\@libraryRepository{\@latex@warning@no@line{No library repository given}}
\def\GetLibraryRepository{\@libraryRepository}

% "(Get|Set)DocumentType".
% The document type of this document.
\def\SetDocumentType#1{\gdef\@documentType{#1}}
\def\@documentType{\@latex@warning@no@line{No document type given}}
\def\GetDocumentType{\@documentType}

\let\tableofcontentsORIG\tableofcontents
\renewcommand\tableofcontents{\tableofcontentsORIG}

% Define "maketitle".
\def\maketitle{%
  \thispagestyle{empty}
  \begin{centering}

  \vspace*{6.0\baselineskip} % White space at the top of the page

  %%%%%%%%%%%%%%%%%%%%%%%%%%%%%%%%
  % TITLE
  %%%%%%%%%%%%%%%%%%%%%%%%%%%%%%%%

  \rule{\textwidth}{1.6pt}\vspace*{-\baselineskip}\vspace*{2pt} % Thick horizontal rule
  \rule{\textwidth}{0.4pt} % Thin horizontal rule

  \vspace{0.75\baselineskip} % Whitespace above the title

  {\normalsize\uppercase{\GetOrganization{}}}\\
  \vspace{1\baselineskip} % For some reason the lines are overlapping.
  {\LARGE\uppercase{\GetLibraryName} \uppercase{\GetLibraryVersion}}

  \vspace{0.75\baselineskip} % Whitespace below the title

  \rule{\textwidth}{0.4pt}\vspace*{-\baselineskip}\vspace{3.2pt} % Thin horizontal rule
  \rule{\textwidth}{1.6pt} % Thick horizontal rule

  \vspace{2\baselineskip} % Whitespace after the title block

  %%%%%%%%%%%%%%%%%%%%%%%%%%%%%%%%%
  % SUBTITLE
  %%%%%%%%%%%%%%%%%%%%%%%%%%%%%%%%%

  \vspace{2\baselineskip} % Whitespace after the title block

  \GetDocumentType % Subtitle or further description

  \vspace*{3\baselineskip} % Whitespace under the subtitle

  %%%%%%%%%%%%%%%%%%%%%%%%%%%%%%%%%
  % EDITORS
  %%%%%%%%%%%%%%%%%%%%%%%%%%%%%%%%%

  \vspace*{1\baselineskip} % Whitespace before editors

  Edited By

  \vspace{0.5\baselineskip} % Whitespace before the editors

  {\scshape\Large \GetAuthor \\
  {\normalsize(\href{mailto:\GetEmail}{\GetEmail})} \\} % Editor list

  \end{centering}

  %%%%%%%%%%%%%%%%%%%%%%%%%%%%%%%%%%
  % Footer
  %%%%%%%%%%%%%%%%%%%%%%%%%%%%%%%%%%
  \vspace*{\fill} % Whitespace before the editors
  \begin{tabularx}{\columnwidth}{m{0.75cm}m{4.75cm}X}
  \multirow{2}{=}{\includegraphics[scale=0.05]{primordialmachine-135x135.png}} & Primordial Machine                                                  & \hfill Created with \LaTeX\\
                                                                               & \href{https://primordialmachine.com}{https://primordialmachine.com} & \hfill Compiled on \today\ at \currenttime
  \end{tabularx}
}

\makeatother

\SetOrganization{Primordial Machine}
\SetLibraryName{Math Matrices}
\SetLibraryVersion{1.2}
\SetLibraryRepository{https://github.com/primordialmachine/math-matrices}
\SetAuthor{Michael Heilmann}
\SetEmail{michaelheilmann@primordialmachine.com}


\SetLibraryIncludeFileName{include.hpp}
\SetLibraryIncludesDirectoryPath{primordialmachine/math-matrices/\newline\$(PlatformTarget.toLower())/\$(Configuration.toLower())/includes}

\SetLibraryIncludeDirectiveFilePath{primordialmachine/math/matrices/include.hpp}

\SetLibraryStaticLibrariesDirectoryPath{primordialmachine/math-matrices/\newline\$(PlatformTarget.toLower())/\$(Configuration.toLower())/libraries}
\SetLibraryStaticLibraryFileName{math-matrices.lib}

\SetDocumentType{User Manual}

\begin{document}

\frontmatter

\begin{titlepage}
\maketitle
\end{titlepage}

\tableofcontents
\addtocontents{toc}{\protect\thispagestyle{empty}}
\pagenumbering{gobble}

\mainmatter

\chapter{Synopsis}
C++ 17 library providing $n \times m$ floating point matrices. 
The library is made available publicly on
\href{\GetLibraryRepository}{Github}
under the
\href{\GetLibraryRepository/blob/master/LICENSE}{MIT License}.

\chapter{Limitations and Restrictions}
The library officially only supports Visual Studio 2017 and Windows 10.

\chapter{Introductory example}
\textit{\color{orange}This library does not provide any examples yet.}
%Examples are located in the \href{\GetLibraryRepository/blob/master/examples}{examples} directory.

%%%%%%%%%%%%%%%%%%%%%%%%%%%%%%%%%%%%%%%%%%%%%%%%%%%%%%%%%%%%%%%%%%%%%%%%%%%%%%%%%%%%%%%%%%%%%%%%%%%
%
% Primordial Machine's Matrices Library
% Copyright (C) 2017-2019 Michael Heilmann
%
% This software is provided 'as-is', without any express or implied warranty.
% In no event will the authors be held liable for any damages arising from the
% use of this software.
%
% Permission is granted to anyone to use this software for any purpose,
% including commercial applications, and to alter it and redistribute it
% freely, subject to the following restrictions:
%
% 1. The origin of this software must not be misrepresented;
%    you must not claim that you wrote the original software.
%    If you use this software in a product, an acknowledgment
%    in the product documentation would be appreciated but is not required.
%
% 2. Altered source versions must be plainly marked as such,
%    and must not be misrepresented as being the original software.
%
% 3. This notice may not be removed or altered from any source distribution.
%
%%%%%%%%%%%%%%%%%%%%%%%%%%%%%%%%%%%%%%%%%%%%%%%%%%%%%%%%%%%%%%%%%%%%%%%%%%%%%%%%%%%%%%%%%%%%%%%%%%%

\chapter{Building under Visual Studio 2017}
\begin{enumerate}
\item Open the solution \texttt{solution.sln} in Microsoft Visual Studio 2017.
\item Batch build everything.
\item The folder \texttt{packages} contains the distribution of the library i.e. include files and the
      static libraries for
  \begin{enumerate}
    \item the platform targets \texttt{x86} and \texttt{x64} and
    \item configurations \texttt{Release} and \texttt{Debug}.
  \end{enumerate}
\item Copy the contents of the \verb+packages+ folder into a directory. Let
      \verb+[library home]+ be a placeholder denoting the path by which that folder
      can be referenced from your project.
\item Add
  \begin{enumerate}
    \item the include path
\texttt{[library home]/\GetLibraryIncludesDirectoryPath}
	and
    \item the library path
\texttt{[library home]/\GetLibraryStaticLibrariesDirectoryPath}
    to your project.
\end{enumerate}
\item Link your project with the library \texttt{\GetLibraryStaticLibraryFileName}.
\item Add the include directive \texttt{\#include "{}\GetLibraryIncludeDirectiveFilePath"{}} where appropriate.
\item You can now use the functionality provided by the library.
\end{enumerate}


\chapter{Library Interface Documentation}

\section{\texttt{namespace primordialmachine}}
The namespace this library is adding its declarations/definitions to.
The added namespace elements are documented below.

%%%%%%%%%%%%%%%%%%%%%%%%%%%%%%%%%%%%%%%%%%%%%%%%%%%%%%%%%%%%%%%%%%%%%%%%%%%%%%%%%%%%%%%%%%%%%%%%%%%
%
% Primordial Machine's Math Points Library
% Copyright (C) 2017-2019 Michael Heilmann
%
% This software is provided 'as-is', without any express or implied warranty.
% In no event will the authors be held liable for any damages arising from the
% use of this software.
%
% Permission is granted to anyone to use this software for any purpose,
% including commercial applications, and to alter it and redistribute it
% freely, subject to the following restrictions:
%
% 1. The origin of this software must not be misrepresented;
%    you must not claim that you wrote the original software.
%    If you use this software in a product, an acknowledgment
%    in the product documentation would be appreciated but is not required.
%
% 2. Altered source versions must be plainly marked as such,
%    and must not be misrepresented as being the original software.
%
% 3. This notice may not be removed or altered from any source distribution.
%
%%%%%%%%%%%%%%%%%%%%%%%%%%%%%%%%%%%%%%%%%%%%%%%%%%%%%%%%%%%%%%%%%%%%%%%%%%%%%%%%%%%%%%%%%%%%%%%%%%%

% Macro to emit the is_* and is_*_v definitions which are used by e.g. angle, color, matrix, point, and vector types.
\newcommand{\Isa}[1]{%
\section{\texttt{is\_#1} (struct template)}
Derives from \texttt{true\_type} is the specified type is a #1 type, derives from \texttt{false\_type} otherwise.%

\subsection{\texttt{is\_#1\_v} (helper variable template)}%
\texttt{template \textlangle typename TYPE\textrangle\newline%
inline constexpr bool is\_#1\_v = is\_#1\textlangle TYPE\textrangle::value;}%
}
\Isa{matrix}

%%%%%%%%%%%%%%%%%%%%%%%%%%%%%%%%%%%%%%%%%%%%%%%%%%%%%%%%%%%%%%%%%%%%%%%%%%%%%%%%%%%%%%%%%%%%%%%%%%%
\section{\texttt{matrix} (struct)}
\begin{verbatim}
template<typename TRAITS, typename ENABLED = void>
struct matrix;
\end{verbatim}
\noindent{}This library provides specializations where \texttt{TRAITS} is a \texttt{matrix\_traits}
specialization with the following properties:
\begin{enumerate}
	\item \texttt{ELEMENT\_TYPE} is a type \texttt{TYPE} for which \texttt{is\_scalar\textlangle TYPE \textrangle::value} is \texttt{true}.
	\item \texttt{NUMBER\_OF\_ROWS} and \texttt{NUMBER\_OF\_COLUMNS} each are of any number between 0 and
	\texttt{std::numeric\_limits\textlangle size\_t\textrangle\::max()}.
\end{enumerate}

\subsection{Members}

\subsection{Arithmetic Functors Support}

%%%%%%%%%%%%%%%%%%%%%%%%%%%%%%%%%%%%%%%%%%%%%%%%%%%%%%%%%%%%%%%%%%%%%%%%%%%%%%%%%%%%%%%%%%%%%%%%%%%
\subsubsection{\texttt{operator+} (binary, matrix matrix)}
This library provides specializations of
\texttt{binary\_plus\_functor}
(see \cite{arithmeticfunctors}) to compute the
sum
of two matrices. Both type arguments must be the same matrix type \texttt{M}, the result type is \texttt{M}.

\subsubsection{\texttt{operator+} (unary, matrix)}
This library provides specializations of
\texttt{unary\_plus\_functor}
(see \cite{arithmeticfunctors}) to compute
the affirmation
of a matrix.
The type argument must be a matrix type \texttt{M}, the result type is \texttt{M}.

\subsubsection{\texttt{operator+=} (compound binary, matrix matrix)}
This library provides specializations of
\texttt{plus\_assignment\_functor}
(see \cite{arithmeticfunctors}) to compute
the sum of to matrices and assign the result to the first matrix.
The type arguments must both the same matrix type \texttt{M}, the result type is \texttt{M}.

%%%%%%%%%%%%%%%%%%%%%%%%%%%%%%%%%%%%%%%%%%%%%%%%%%%%%%%%%%%%%%%%%%%%%%%%%%%%%%%%%%%%%%%%%%%%%%%%%%%
\subsubsection{\texttt{operator-} (binary, matrix matrix)}
This library provides specializations of
\texttt{binary\_minus\_functor}
(see \cite{arithmeticfunctors}) to compute the
difference
of two matrices. Both type arguments must be the same matrix type \texttt{M}, the result type is \texttt{M}.

\subsubsection{\texttt{operator-} (unary, matrix)}
This library provides specializations of
\texttt{unary\_minus\_functor}
(see \cite{arithmeticfunctors}) to compute
the negation
of a matrix.
The type argument must be a matrix type \texttt{M}, the result type is \texttt{M}.

\subsubsection{\texttt{operator-=} (compound binary, matrix matrix)}
This library provides specializations of
\texttt{minus\_assignment\_functor}
(see \cite{arithmeticfunctors}) to compute
the difference of to matrices and assign the result to the first matrix.
The type arguments must both the same matrix type \texttt{M}, the result type is \texttt{M}.

%%%%%%%%%%%%%%%%%%%%%%%%%%%%%%%%%%%%%%%%%%%%%%%%%%%%%%%%%%%%%%%%%%%%%%%%%%%%%%%%%%%%%%%%%%%%%%%%%%%
\subsubsection{\texttt{operator*} (binary, matrix scalar)}
This library provides specializations of
\texttt{binary\_star\_functor}
(see \cite{arithmeticfunctors} to compute the
product of a matrix and a scalar.
The first type argument is a matrix type \texttt{M}, the second type parameter is a scalar type \texttt{S} (see \cite{mathscalars}
such that \texttt{is\_same\_v\textlangle element\_type\_v\textlangle M\textrangle, S\textrangle} holds.
The result type is \texttt{M}.

\subsubsection{\texttt{operator*=} (compound binary, matrix scalar)}
This library provides specializations of
\texttt{star\_assignment\_functor}
(see \cite{arithmeticfunctors}) to compute the
product of a matrix and a scalar and assign the result to the matrix.
The first type argument is a matrix type \texttt{M}, the second type parameter is a scalar type \texttt{S} (see \cite{mathscalars}
such that \texttt{is\_same\_v\textlangle element\_type\_v\textlangle M\textrangle, S\textrangle} holds.
The result type is \texttt{M}.


\subsubsection{\texttt{operator*} (binary, scalar matrix)}
This library provides specializations of
\texttt{binary\_star\_functor}
(see \cite{arithmeticfunctors} to compute the
product of a scalar and a matrix.
The first type parameter is a scalar type \texttt{S} (see \cite{mathscalars}), the second type argument is a matrix type \texttt{M}.
such that \texttt{is\_same\_v\textlangle element\_type\_v\textlangle M\textrangle, S\textrangle} holds.
The result type is \texttt{M}.

\subsubsection{\texttt{operator/} (binary, matrix scalar)}
This library provides specializations of
\texttt{binary\_star\_functor}
(see \cite{arithmeticfunctors} to compute the
product of a matrix and a scalar.
The first type argument is a matrix type \texttt{M}, the second type parameter is a scalar type \texttt{S} (see \cite{mathscalars}
such that \texttt{is\_same\_v\textlangle element\_type\_v\textlangle M\textrangle, S\textrangle} holds.
The result type is \texttt{M}.

\subsubsection{\texttt{operator/=} (compound binary, matrix scalar)}
This library provides specializations of
\texttt{star\_assignment\_functor}
(see \cite{arithmeticfunctors}) to compute the
quotient of a matrix and a scalar and assign the result to the matrix.
The first type argument is a matrix type \texttt{M}, the second type parameter is a scalar type \texttt{S} (see \cite{mathscalars}
such that \texttt{is\_same\_v\textlangle element\_type\_v\textlangle M\textrangle, S\textrangle} holds.
The result type is \texttt{M}.

\subsubsection{\texttt{operator*} (binary, matrix matrix)}
This library provides specializations of
\texttt{binary\_star\_functor}
(see \cite{arithmeticfunctors}) to compute the
product
of two matrices. The first type argument is a matrix type \texttt{A} and the second type argument is a matrix type \texttt{B}
such that
\begin{enumerate}
	\item \texttt{number\_of\_columns\_v\textlangle A \textrangle} equals \texttt{number\_of\_rows\_v\textlangle B\textrangle}.
	\item \texttt{is\_same\_v\textlangle element\_type\_v\textlangle A, B\textrangle\textrangle}
\end{enumerate}
The result type is a matrix type \texttt{C} such that
\begin{enumerate}
	\item \texttt{number\_of\_rows\_v\textlangle C\textrangle := number\_of\_rows\_v\textlangle A\textrangle}
	\item \texttt{number\_of\_columns\_v\textlangle C\textrangle := number\_of\_columns\_v\textlangle B\textrangle}
	\item \texttt{element\_type\_v\textlangle C\textrangle := element\_type\_v\textlangle A\textrangle}
\end{enumerate}

%%%%%%%%%%%%%%%%%%%%%%%%%%%%%%%%%%%%%%%%%%%%%%%%%%%%%%%%%%%%%%%%%%%%%%%%%%%%%%%%%%%%%%%%%%%%%%%%%%%

\subsection{Relational Functors Support}

\subsubsection{\texttt{operator==} (compound binary, matrix matrix)}
This library provides specializations of
\texttt{equal\_to\_functor}
(see \cite{relationalfunctors}) to compute
wether two matrices of the same type are equal.

\subsubsection{\texttt{operator!=} (compound binary, matrix matrix)}
This library provides specializations of
\texttt{not\_equal\_to\_functor}
(see \cite{relationalfunctors}) to compute
wether two matrices of the same type are not equal.

%%%%%%%%%%%%%%%%%%%%%%%%%%%%%%%%%%%%%%%%%%%%%%%%%%%%%%%%%%%%%%%%%%%%%%%%%%%%%%%%%%%%%%%%%%%%%%%%%%%

\subsection{\texttt{inverse\_functor} (\textit{FunctorBase})}
A \textit{FunctorBase} with a single template parameter \texttt{M} followed by the \texttt{ENABLED} template parameter.
Specializations of this \textit{FunctorBase} compute inverses of matrices.
A possible implementation is given by
\begin{verbatim}
template<typename A, typename ENABLED = void>
struct transpose_functor;
\end{verbatim}

\subsubsection{\texttt{inverse} (function)}
This library provides a function to invoke the \texttt{inverse\_functor}.
A possible implementation is given by\newline
\texttt{template\textlangle typename A\textrangle\newline
auto\newline
inverse(const A\& a)\newline
noexcept(noexcept(inverse\_functor\textlangle A \textrangle))\newline
-> decltype(inverse\_functor\textlangle A\textrangle()(operand))\newline
\{ return inverse\_functor\textlangle A\textrangle()(a); \}}

\subsection{\texttt{inverse\_functor} (\textit{Functor})}
This library provides \textit{Functor} specializations of the \texttt{inverse\_functor} \textit{FunctorBase}
where \texttt{A} is a non degenerate square matrix type and the result type is \texttt{A}.

\subsection{\texttt{transpose\_functor} (\textit{FunctorBase})}
A \textit{FunctorBase} with a single parameter \texttt{A} followed by the \texttt{ENABLED} template parameter.
Specializations of this \textit{FunctorBase} compute transposes of matrices.
A possible implementation is given by
\begin{verbatim}
template<typename A, typename ENABLED = void>
struct transpose_functor;
\end{verbatim}

\subsubsection{\texttt{transpose} (function)}
This library provides a function to invoke the \texttt{transpose\_functor}.
A possible implementation is given by\newline
\texttt{template\textlangle typename A\textrangle\newline
auto\newline
transpose(const A\& a)\newline
noexcept(noexcept(transpose\_functor\textlangle A\textrangle()(a)))\newline
-> decltype(transpose\_functor\textlangle A\textrangle()(a))\newline
\{ return transpose\_functor\textlangle A\textrangle()(a); \}}

\subsection{\texttt{transpose\_functor} (\textit{Functor})}
This library provides \textit{Functor} specializations of the \texttt{transpose\_functor} \textit{FunctorBase}
where \texttt{A} is matrix type and the result type is \texttt{B} such that
\begin{enumerate}
	\item \texttt{element\_type\_v\textlangle B\textrangle = element\_type\_v\textlangle A\textrangle}
	\item \texttt{number\_of\_rows\_v\textlangle B \textrangle = number\_of\_columns\_v\textlangle A\textrangle}
	\item \texttt{number\_of\_columns\_v\textlangle B \textrangle = number\_of\_rows\_v\textlangle A\textrangle}
\end{enumerate}
\noindent{}Note that the number of rows and columns of the result type are swapped.

\subsection{\texttt{trace\_functor} (\textit{FunctorBase})}
An \textit{FunctorBase} with a single template parameter \texttt{A} followed by the \texttt{ENALBED} parameter.
Specializations of this \textit{FunctorBase} compute traces of matrices.
A possible implementation is given by
\begin{verbatim}
template<typename A, typename ENABLED = void>
struct trace_functor;
\end{verbatim}

\subsubsection{\texttt{trace} (function)}
This library provides a function to invoke the \texttt{trace\_functor}.
A possible implementation is given by\newline
\texttt{template\textlangle typename A\textrangle\newline
auto\newline
trace(const A\& a)\newline
noexcept(noexcept(trace\_functor\textlangle A\textrangle()(a)))
-> decltype(trace\_functor\textlangle A\textrangle()(a))\newline
\{ return trace\_functor\textlangle A\textrangle()(a); \}}

\subsection{\texttt{trace\_functor} (\textit{Functor})}
This library provides \textit{Functor} specializations of the \texttt{trace\_functor} \textit{FunctorBase}
where \texttt{A} is a non degenerate matrix type and the result type is \texttt{element\_type\_v\textlangle
A\textrangle}.


%%%%%%%%%%%%%%%%%%%%%%%%%%%%%%%%%%%%%%%%%%%%%%%%%%%%%%%%%%%%%%%%%%%%%%%%%%%%%%%%%%%%%%%%%%%%%%%%%%%
%
% Primordial Machine Math Matrices Library
% Copyright (c) 2017-2019 Michael Heilmann <michaelheilmann@primordialmachine.com>
%
% This software is provided 'as-is', without any express or implied warranty.
% In no event will the authors be held liable for any damages arising from the
% use of this software.
%
% Permission is granted to anyone to use this software for any purpose,
% including commercial applications, and to alter it and redistribute it
% freely, subject to the following restrictions:
%
% 1. The origin of this software must not be misrepresented;
%    you must not claim that you wrote the original software.
%    If you use this software in a product, an acknowledgment
%    in the product documentation would be appreciated but is not required.
%
% 2. Altered source versions must be plainly marked as such,
%    and must not be misrepresented as being the original software.
%
% 3. This notice may not be removed or altered from any source distribution.
%
%%%%%%%%%%%%%%%%%%%%%%%%%%%%%%%%%%%%%%%%%%%%%%%%%%%%%%%%%%%%%%%%%%%%%%%%%%%%%%%%%%%%%%%%%%%%%%%%%%%

\nocite{*} % Force all bibliography items to be printed, regardless of wether they are cited or not.

\begin{thebibliography}{9}
\bibitem{functors}
Primordial Machine's Functors Library 
\\\texttt{https://github.com/primordialmachine/functors}

\bibitem{arithmeticfunctors}
Primordial Machine's Arithmetic Functors Library 
\\\texttt{https://github.com/primordialmachine/arithmetic-functors}

\bibitem{relationalfunctors}
Primordial Machine's Relational Functors Library 
\\\texttt{https://github.com/primordialmachine/relational-functors}

\bibitem{mathscalars}
Primordial Machine's Math Scalars Library 
\\\texttt{https://github.com/primordialmachine/math-scalars}

\bibitem{mathvectors}
Primordial Machine's Math Vectors Library 
\\\texttt{https://github.com/primordialmachine/math-vectors}

\bibitem{mathmatrices}
Primordial Machine's Math Matrices Library 
\\\texttt{https://github.com/primordialmachine/math-matrices}


\end{thebibliography}


\end{document}
